% ==================================================================
% Hawking Radiation and the Planck Remnant
% in Conscious Point Physics
% Companion 10 to "Mechanistic Derivation of Relativistic Effects
% via Space Stress Vector (SSV) in the Dipole Sea" (Version 16)
% Version 1 — 18 March 2026
% ==================================================================

\documentclass[12pt]{article}

\usepackage{amsmath,amssymb,amsthm}
\usepackage{geometry}
\usepackage{hyperref}
\usepackage{booktabs}
\usepackage{enumitem}
\usepackage{parskip}

\geometry{letterpaper, margin=1in}

\newtheorem{theorem}{Theorem}[section]
\newtheorem{proposition}[theorem]{Proposition}
\newtheorem{corollary}[theorem]{Corollary}
\newtheorem{remark}[theorem]{Remark}

% ── CPP macros ────────────────────────────────────────────────────────────────
\newcommand{\lP}{l_P}
\newcommand{\tP}{t_P}
\newcommand{\mP}{m_P}
\newcommand{\EP}{E_P}
\newcommand{\TP}{T_P}
\newcommand{\rS}{r_S}
\newcommand{\rcore}{r_{\mathrm{core}}}
\newcommand{\PSR}{\mathrm{PSR}}
\newcommand{\SSV}{\mathrm{SSV}}
\newcommand{\TH}{T_H}

\title{\textbf{Hawking Radiation and the Planck Remnant\\
in Conscious Point Physics}\\[6pt]
{\large Companion~10 to ``Mechanistic Derivation of Relativistic Effects\\
via Space Stress Vector (SSV) in the Dipole Sea'' (Version~16)}}

\author{Thomas Lee Abshier, ND\\
Hyperphysics Institute\\
\texttt{https://hyperphysics.com}\\
\texttt{drthomas007@protonmail.com}}

\date{18 March 2026 --- Version~1}

\begin{document}
\maketitle

\noindent\textbf{Keywords:} Conscious Point Physics, Hawking radiation,
Planck remnant, black hole evaporation, CP Exclusion, geometric breakdown,
information paradox, black hole thermodynamics.

%======================================================================
\section*{Plain Language Summary}
%======================================================================

Black holes slowly radiate energy — a process discovered theoretically
by Stephen Hawking in 1974.  In general relativity, this radiation
continues until the black hole completely evaporates, ending in a
singular explosion.  The Planck core established in companion~8 changes
this ending.

In CPP, evaporation proceeds normally until the black hole reaches the
mass of a single Planck unit ($m_P \approx 22\,\mu$g, about the mass
of a grain of sand).  At this point, the Schwarzschild radius is only
about twice the Planck lattice spacing.  The Hawking calculation requires
a smooth horizon region much larger than the lattice spacing — a condition
that fails completely when only two lattice cells separate the horizon
from the core.  Evaporation terminates, not because the temperature
gets too high, but because the geometry that allows thermal radiation
no longer exists at the Planck scale.

What remains is a Planck-mass ``remnant'': a configuration of Conscious
Points compressed to the maximum packing density, held stable by
an exact mechanical balance between gravitational SSV compression and
CP Exclusion repulsion.  This remnant does not need external support.
It is self-sustaining because the CP Exclusion bouncing that prevents
further compression is also what keeps the CPs from dispersing.
The remnant interacts gravitationally with its surroundings but cannot
radiate further.

%======================================================================
\begin{abstract}
We derive the CPP modifications to Hawking black hole evaporation
resulting from the Planck core established in companion~8.
Standard Hawking radiation proceeds unchanged for $M \gg \mP$.
Evaporation terminates when the Schwarzschild radius reaches
$\rS \sim 2\lP$ — not when the Hawking temperature reaches the Planck
temperature, but when the geometric condition $\rS \gg \lP$ required
by the semiclassical approximation breaks down.
The result is a stable Planck-mass remnant $M_{\rm rem} \sim \mP$
(Proposition~\ref{prop:remnant}).  We show that this remnant is
mechanically stable via a self-sustaining fixed point: gravitational
SSV compression is exactly balanced by CP Exclusion repulsion at
$\PSR_{\rm eff} = \lP/2$ (Proposition~\ref{prop:stability}).
The remnant does not require external pressure and cannot evaporate
further.  We discuss implications for the black hole information paradox
and observable consequences.  No new postulates beyond companion~8 are
required.
\end{abstract}

%======================================================================
\section{Introduction}
\label{sec:intro}
%======================================================================

Hawking~\cite{hawking1975} showed that black holes are not perfectly
black: quantum effects near the horizon cause them to emit thermal
radiation at temperature $\TH = \hbar c^3/(8\pi G M k_B)$.  This
radiation carries energy away from the black hole, reducing its mass.
In GR, the process continues indefinitely as $M \to 0$, with the
temperature diverging and the evaporation rate accelerating without
bound — ending in a Planck-scale explosion whose final state is
theoretically unresolved.

The CPP strong-field companion (companion~8)~\cite{c8} established
that the CP Exclusion Rule prevents $\PSR_{\rm eff} < \lP/2$,
replacing the GR singularity with a Planck-density core at
$\rcore = \rS/2$.  This inner boundary modifies the Hawking calculation
in two ways.  First, as established in the GW echo companion
(companion~9)~\cite{c9}, the Planck core acts as a reflective surface
whose presence modifies the boundary conditions for quantum field
propagation in the black hole spacetime.  Second, and more
fundamentally, the Planck core sets a minimum black hole mass below
which the geometric conditions for Hawking radiation cannot be met.

The present companion derives both modifications and identifies the
stable Planck remnant as the endpoint of CPP black hole evaporation.
The derivation uses thermodynamic and geometric arguments (Option~B)
rather than the full Bogoliubov coefficient calculation (deferred as
Open Problem~\ref{op:bogoliubov}).

%======================================================================
\section{Standard Hawking Radiation: A Review}
\label{sec:hawking_review}
%======================================================================

The standard Hawking temperature for a Schwarzschild black hole of
mass $M$ is~\cite{hawking1975,wald1984}:
\begin{equation}
  \TH(M) \;=\; \frac{\hbar c^3}{8\pi G M k_B}
  \;=\; \frac{\TP}{8\pi}\cdot\frac{\mP}{M},
  \label{eq:TH}
\end{equation}
where $\TP = \mP c^2/k_B = \sqrt{\hbar c^5/G}/k_B \approx 1.42\times 10^{32}$~K
is the Planck temperature.  The radiation has a thermal (Planck) spectrum
and carries away energy at the Stefan-Boltzmann luminosity:
\begin{equation}
  L_H(M) \;=\; \frac{\hbar c^6}{15360\pi G^2 M^2}.
  \label{eq:LH}
\end{equation}
The mass loss rate $dM/dt = -L_H/c^2$ integrates to the standard
evaporation time:
\begin{equation}
  t_{\rm evap}(M_0) \;=\; \frac{5120\pi G^2 M_0^3}{\hbar c^4}.
  \label{eq:tevap}
\end{equation}
For a solar-mass black hole: $t_{\rm evap}(M_\odot) \approx 2.1\times10^{67}$~yr
— vastly longer than the current age of the universe.

These results hold in the \emph{semiclassical approximation}, which
requires the horizon region to be smooth on scales much larger than
the Planck length:
\begin{equation}
  N \;\equiv\; \frac{\rS}{\lP} \;=\; \frac{2GM}{c^2\lP}
  \;=\; 2\,\frac{M}{\mP} \;\gg\; 1.
  \label{eq:N}
\end{equation}
This is the condition that the quantum field modes relevant to Hawking
radiation see a slowly-varying background geometry.

%======================================================================
\section{The Geometric Termination Condition}
\label{sec:termination}
%======================================================================

\begin{proposition}[Geometric termination]
\label{prop:termination}
Hawking evaporation terminates when $M = M_{\rm rem} \sim \mP$,
where the semiclassical approximation breaks down geometrically.
\end{proposition}

\begin{proof}
The semiclassical approximation requires $N = \rS/\lP \gg 1$,
i.e.\ $M \gg \mP/2$.  As $M$ decreases toward $\mP$:
\begin{equation}
  N(M) \;=\; \frac{2M}{\mP} \;\longrightarrow\; 2
  \quad\text{as } M \to \mP.
  \label{eq:N_mP}
\end{equation}
At $M = \mP$, the Schwarzschild radius is:
\begin{equation}
  \rS(\mP) \;=\; \frac{2G\mP}{c^2} \;=\; 2\lP,
  \label{eq:rS_mP}
\end{equation}
spanning only two Planck lattice cells.  The Hawking mechanism
requires a smooth thermal gradient across the near-horizon region.
In the 600-cell lattice, there is no such region when $\rS \sim 2\lP$:
the ``horizon'' is a boundary of width $\lP$ separating two lattice
cells from three.  The semiclassical approximation is not merely
inaccurate — it is geometrically undefined.

Therefore, the Hawking calculation is valid for $M \gg \mP$ and
terminates at $M \sim \mP$.
\end{proof}

\begin{remark}[Temperature at termination]
One might expect termination when $\TH \to \TP$ (the Planck
temperature).  This does not occur first.  At $M = \mP$:
\begin{equation}
  \TH(\mP) \;=\; \frac{\TP}{8\pi} \;\approx\; 0.040\,\TP.
  \label{eq:TH_mP}
\end{equation}
The Hawking temperature at termination is only $4\%$ of the Planck
temperature — well below the lattice maximum.  The geometric
breakdown ($N = 2$) is the binding constraint, not the thermal one.
\end{remark}

\begin{remark}[Evaporation time to the remnant]
The total evaporation time from initial mass $M_0$ to the Planck
remnant $\mP$ is:
\begin{equation}
  t_{\rm evap}(M_0 \to \mP)
  \;\approx\; \frac{5120\pi G^2 M_0^3}{\hbar c^4}
  \left[1 - \left(\frac{\mP}{M_0}\right)^3\right]
  \;\approx\; t_{\rm evap}(M_0),
  \label{eq:tevap_remnant}
\end{equation}
since $(\mP/M_0)^3 \ll 1$ for any astrophysical black hole.
The total evaporation time is essentially unchanged from the
standard result: $\sim\!2.1\times10^{67}$~yr for a solar-mass
black hole.
\end{remark}

%======================================================================
\section{The Planck Remnant}
\label{sec:remnant}
%======================================================================

\begin{proposition}[Planck remnant]
\label{prop:remnant}
The endpoint of CPP black hole evaporation is a stable Planck-mass
remnant with:
\begin{align}
  M_{\rm rem} &\;\sim\; \mP \;=\; \sqrt{\hbar c/G}
  \;\approx\; 2.18\times 10^{-8}\ \mathrm{kg}, \label{eq:Mrem}\\
  r_{\rm rem} &\;=\; \frac{G\mP}{c^2} \;=\; \lP
  \;\approx\; 1.62\times 10^{-35}\ \mathrm{m}, \label{eq:rrem}\\
  \rho_{\rm rem} &\;\sim\; \frac{\mP}{\lP^3}
  \;=\; \frac{c^5}{\hbar G^2}
  \;\approx\; 5.16\times 10^{96}\ \mathrm{kg\,m^{-3}}. \label{eq:rhoRem}
\end{align}
\end{proposition}

\begin{proof}
From Proposition~\ref{prop:termination}, Hawking evaporation stops at
$M \sim \mP$.  The CP Exclusion Rule~\cite{c1} prevents further
compression: $\PSR_{\rm eff} \geq \lP/2$ enforces $r \geq \rcore = GM/c^2$.
At $M = \mP$, $\rcore = \lP$.  No further mass loss is possible
because the Hawking mechanism has terminated, and no further
compression is possible because the CP Exclusion floor has been
reached.  The remnant properties follow directly from $M = \mP$,
$r = \lP$ in Eq.~\eqref{eq:rhoRem}.
\end{proof}

Table~\ref{tab:remnant} summarises the remnant properties.

\begin{table}[h]
\centering
\caption{Properties of the CPP Planck remnant.}
\label{tab:remnant}
\renewcommand{\arraystretch}{1.35}
\begin{tabular}{@{}lll@{}}
\toprule
Quantity & Value & Expression \\
\midrule
Remnant mass $M_{\rm rem}$
  & $2.176\times 10^{-8}$~kg
  & $\sqrt{\hbar c/G}$ \\
Remnant radius $r_{\rm rem}$
  & $1.616\times 10^{-35}$~m
  & $\lP$ \\
Remnant density $\rho_{\rm rem}$
  & $5.16\times 10^{96}$~kg\,m$^{-3}$
  & $\mP/\lP^3 = c^5/(\hbar G^2)$ \\
$N = \rS/\lP$ at termination
  & 2 & $2\mP/\mP$ \\
$T_H$ at termination
  & $5.64\times 10^{30}$~K $= \TP/8\pi$
  & $\TP\mP/(8\pi M_{\rm rem})$ \\
Fraction of mass radiated
  & $\approx 100\%$
  & $1 - \mP/M_0 \approx 1$ for $M_0 \gg \mP$ \\
\bottomrule
\end{tabular}
\end{table}

%======================================================================
\section{Stability of the Planck Remnant}
\label{sec:stability}
%======================================================================

\begin{proposition}[Mechanical stability of the remnant]
\label{prop:stability}
The Planck remnant is mechanically stable without external pressure.
It is sustained by an exact balance between gravitational SSV
compression and CP Exclusion repulsion at the PSR saturation point.
\end{proposition}

\begin{proof}
The CP Exclusion Rule (Absolute Moment companion~\cite{c1}) operates
as follows: when two Conscious Points arrive at the same Grid Point
address on the same Absolute Moment tick, they do not stack.  On the
immediately following tick, each CP moves outward along the PSR shell
in the direction defined by the local $\SSV_{\rm net}$ vector — which
points radially outward from the compressed core.

Consider the PSR formula at the core:
\begin{equation}
  \PSR_{\rm eff}(r) \;=\; \frac{\lP}{1 + GM/rc^2}.
  \label{eq:psr_core}
\end{equation}
The gravitational SSV compression force drives $\PSR_{\rm eff}$ downward
(higher compression).  The CP Exclusion repulsion force drives CPs
outward when $\PSR_{\rm eff}$ approaches $\lP/2$.  These two forces
are exactly balanced at the saturation point:
\begin{equation}
  \PSR_{\rm eff}(\rcore) \;=\; \frac{\lP}{2}
  \quad\Longleftrightarrow\quad
  \frac{GM}{\rcore c^2} \;=\; 1
  \quad\Longleftrightarrow\quad
  \rcore \;=\; \frac{GM}{c^2}.
  \label{eq:saturation}
\end{equation}
This is a \emph{mechanical fixed point}: any perturbation that
compresses the remnant below $\rcore$ is immediately countered by
increased CP Exclusion repulsion; any perturbation that expands it
above $\rcore$ is countered by the increase in gravitational SSV.
The balance is not tuned — it is an exact structural consequence of
the PSR formula and the CP Exclusion floor.

The remnant therefore requires no external pressure.  The CPs
constituting the remnant are sustained by mutual bouncing:
each CP's outward impulse from CP Exclusion defines the SSV field
experienced by its neighbours, which defines their exclusion
directions in the next tick.  This is a self-referential,
self-sustaining dynamic equilibrium.
\end{proof}

\begin{remark}[Why the remnant cannot evaporate further]
Hawking radiation requires a thermal gradient across a smooth
near-horizon region.  This requires $\rS \gg \lP$.  At the remnant
mass $M = \mP$, the Schwarzschild radius is $\rS = 2\lP$ — only two
lattice spacings.  There is no smooth thermal gradient; the entire
structure is at the Planck scale.  The Hawking mechanism has no
meaning for a two-cell object.  What remains is a purely mechanical
system of CPs executing the PCD cycle at maximum compression —
not a thermally radiating black hole.
\end{remark}

\begin{remark}[Comparison with neutron star stability]
A neutron star is stabilised by neutron degeneracy pressure, which
requires a macroscopic quantum Fermi sea — an external environment of
neutrons.  Remove the neutrons and the star collapses.  The CPP Planck
remnant is fundamentally different: it requires no external medium.
The CPs \emph{are} the pressure; their mutual CP Exclusion bouncing
\emph{is} the stabilising force.  The remnant is stable in vacuum,
in the absence of any surrounding matter.
\end{remark}

%======================================================================
\section{The Information Paradox in CPP}
\label{sec:information}
%======================================================================

The black hole information paradox asks: where does the quantum
information encoded in the infalling matter go?  In GR, the singularity
destroys it.  In CPP, there is no singularity, so this destruction
does not occur.  Instead:

\begin{itemize}[noitemsep]
  \item The Hawking radiation during the evaporation phase carries
    thermal (mixed-state) radiation, consistent with semiclassical
    expectations.

  \item The Planck remnant retains the quantum information encoded in
    its internal CP configuration — specifically, the pattern of
    $\SSV_{\rm net}$ vectors and CP positions within the $\sim 2\lP$
    core.  The information capacity of a Planck-volume core is
    $S \sim k_B \ln 2$ per CP, or of order $k_B$ total — the
    minimum possible entropy consistent with the Bekenstein bound
    at Planck scale.

  \item Whether full unitarity is preserved — i.e.\ whether the
    combined Hawking radiation plus remnant contains all the
    original information — is the CPP analog of the information
    paradox.  This requires the full quantum treatment
    (Open Problem~\ref{op:bogoliubov}) and is deferred.
\end{itemize}

\begin{remark}[Page time and the remnant]
Page~\cite{page1993} showed that if Hawking radiation is unitary, the
radiation must begin to encode information after the ``Page time''
$t_{\rm Page} \approx t_{\rm evap}/2$.  In CPP, the modification
is at the very end of evaporation ($M \to \mP$), which is a negligibly
small fraction of $t_{\rm evap}$.  The early-time Hawking radiation is
unchanged; the information question is deferred to the final Planck
epoch.
\end{remark}

%======================================================================
\section{Consistency with Previous Companions}
\label{sec:consistency}
%======================================================================

\begin{itemize}[noitemsep]
  \item The \emph{CP Exclusion Rule} establishing $\PSR_{\rm eff} \geq \lP/2$
    is from companion~1~\cite{c1} and companion~8~\cite{c8}.
  \item The \emph{Planck core} at $\rcore = \rS/2$ is from companion~8,
    Theorem~3.
  \item $G = \hbar c/\mP^2$ giving $\rS/\lP = 2M/\mP$ is from
    companion~5~\cite{c5}.
  \item The \emph{GW echo} signal from the same Planck core is
    companion~9~\cite{c9}; the two papers share the same inner
    boundary and cite each other.
  \item The \emph{exact Schwarzschild exterior} (companion~8,
    Theorem~2) ensures the Hawking calculation is unchanged for
    $M \gg \mP$: the CPP exterior is identical to GR at all
    $r > \rS$.
\end{itemize}

No new postulates beyond companion~8 are required.

%======================================================================
\section{Open Problems}
\label{sec:open}
%======================================================================

\begin{enumerate}
  \item \label{op:bogoliubov}
    \textbf{Full Bogoliubov coefficient calculation.}
    The thermodynamic argument of Proposition~\ref{prop:termination}
    identifies the termination condition geometrically.  A complete
    derivation would compute the Bogoliubov coefficients for quantum
    field modes propagating between the asymptotic past and the
    near-Planck epoch, with the inner boundary condition at
    $r_0 = \rS + \lP$ (CP Exclusion floor).  This is the CPP analog
    of the full Hawking derivation and would determine whether the
    late-time spectrum departs from thermal before termination.

  \item \textbf{Information recovery.}
    Does the CPP remnant carry sufficient information to resolve the
    information paradox?  The information capacity of a Planck-volume
    core is of order $k_B$, which is far less than the
    Bekenstein--Hawking entropy $S_{\rm BH} = k_B A/(4\lP^2)$ of
    the original black hole.  The resolution must involve the full
    Hawking radiation train, not just the remnant.

  \item \textbf{Remnant interactions.}
    The Planck remnant interacts gravitationally ($G = \hbar c/\mP^2$,
    companion~5) but its interactions via charge-polarity SSV
    (electromagnetism) depend on whether the residual CP charge
    configuration inside the remnant is net neutral.  The charge
    state of the remnant after evaporation is an open question.

  \item \textbf{Cosmological remnant density.}
    If primordial black holes formed in the early universe and have
    evaporated to Planck remnants, the cosmological remnant density
    provides a bound on primordial black hole abundance.  Deriving
    this bound within CPP is deferred.
\end{enumerate}

%======================================================================
\section{Conclusion}
\label{sec:conclusion}
%======================================================================

\begin{enumerate}
  \item \textbf{Hawking radiation unchanged for $M \gg \mP$}
    (standard result preserved): the CPP exterior is identical to GR
    for $M \gg \mP$; the standard Hawking luminosity and evaporation
    time apply without modification.

  \item \textbf{Geometric termination at $M \sim \mP$}
    (rigorous, Proposition~\ref{prop:termination}): the
    semiclassical approximation requires $\rS \gg \lP$, equivalently
    $N = 2M/\mP \gg 1$.  At $M = \mP$, $N = 2$ and the geometric
    condition fails.  Evaporation terminates at the remnant mass,
    not at a temperature divergence: $\TH(\mP) = \TP/8\pi \approx 0.04\,\TP$.

  \item \textbf{Planck remnant mass and density}
    (rigorous, Proposition~\ref{prop:remnant}):
    $M_{\rm rem} \sim \mP \approx 2.18\times10^{-8}$~kg,
    $r_{\rm rem} = \lP$,
    $\rho_{\rm rem} \approx 5.16\times10^{96}$~kg\,m$^{-3}$.

  \item \textbf{Remnant stability without external pressure}
    (rigorous, Proposition~\ref{prop:stability}): the CP Exclusion
    Rule creates a mechanical fixed point at $\PSR_{\rm eff} = \lP/2$
    that exactly balances gravitational SSV compression.  The remnant
    is self-sustaining via internal CP bouncing dynamics; no external
    environment is required.

  \item \textbf{No further evaporation} (rigorous): the Hawking
    mechanism requires a smooth near-horizon geometry ($\rS \gg \lP$)
    that does not exist for the two-cell Planck remnant.  What remains
    is a purely mechanical fixed-point system, not a thermally
    radiating object.

  \item \textbf{Information paradox deferred}: the Planck remnant
    retains order-$k_B$ information, far less than the original
    Bekenstein--Hawking entropy.  Full unitarity requires analysis
    of the complete Hawking radiation train (Open
    Problem~\ref{op:bogoliubov}).
\end{enumerate}

The Planck remnant completes the CPP black hole picture: exact GR
exterior, echoing Planck-core interior (companion~9), and stable
Planck-mass endpoint.  Three papers — companions 8, 9, and 10 —
together form a complete, falsifiable, internally consistent account
of the CPP black hole from formation to final state.

%======================================================================
\section*{Acknowledgments}
%======================================================================

The identification of the geometric breakdown ($N = 2$ at $M = \mP$)
as the correct termination condition — rather than the temperature
divergence — and the mechanical fixed-point argument for remnant
stability were established on 18 March 2026.  The author thanks
Grok (xAI) and Claude (Anthropic) for contributions to the CPP
black hole series.

%======================================================================
\begin{thebibliography}{99}

\bibitem{v16}
T.~L. Abshier,
``Mechanistic Derivation of Relativistic Effects via Space Stress
Vector (SSV) in the Dipole Sea,''
Version~16, 2026 (main paper).

\bibitem{c1}
T.~L. Abshier,
``The Absolute Moment Postulate,''
Version~3, 2026 (companion~1).

\bibitem{c5}
T.~L. Abshier,
``Newtonian Gravity from SSV Shell Broadcast,''
Version~2, 2026 (companion~5).

\bibitem{c8}
T.~L. Abshier,
``Strong-Field General Relativity and the CPP Field Equation,''
Version~2, 2026 (companion~8).

\bibitem{c9}
T.~L. Abshier,
``Gravitational Wave Echoes from the Planck Core,''
Version~2, 2026 (companion~9).

\bibitem{hawking1975}
S.~W. Hawking,
``Particle Creation by Black Holes,''
\textit{Commun.\ Math.\ Phys.}, \textbf{43}, 199--220, 1975.

\bibitem{bekenstein1973}
J.~D. Bekenstein,
``Black Holes and Entropy,''
\textit{Phys.\ Rev.\ D}, \textbf{7}, 2333--2346, 1973.

\bibitem{page1993}
D.~N. Page,
``Information in Black Hole Radiation,''
\textit{Phys.\ Rev.\ Lett.}, \textbf{71}, 3743--3746, 1993.

\bibitem{wald1984}
R.~M. Wald,
\textit{General Relativity},
University of Chicago Press, Chicago, 1984.

\bibitem{adler2001}
R.~J. Adler, P.~Chen, and D.~I. Santiago,
``The Generalized Uncertainty Principle and Black Hole Remnants,''
\textit{Gen.\ Rel.\ Grav.}, \textbf{33}, 2101--2108, 2001.

\bibitem{codata}
NIST, CODATA 2018 Recommended Values of the Fundamental Physical
Constants, \url{https://physics.nist.gov/cuu/Constants/}, 2018.

\end{thebibliography}

\end{document}
